\section{Local factors of the chronological source expansion}
\label{sec:peps_chronological_local_factors}

\begin{definition}[Labels of participating occurrences]
    \label{def:peps_nonempty_gate_labels}
    \lean{TNLean.PEPS.PairEffect.SourceCircuit.nonemptyGateLocations,
      TNLean.PEPS.PairEffect.SourceCircuit.allChoicesNonemptyLabelsEquiv}
    \leanok
    Let $W$ be a source circuit, and let $\mathcal G_W$ be its ordered
    gate occurrences. Write $C_g$ for the participants of occurrence $g$
    and $\Xi_g$ for its monomial labels. Set
    $\mathcal G_W^+=\{g\in\mathcal G_W:C_g\ne\varnothing\}$.
    Choices in the expansion over all parties are in bijection with
    $\prod_{g\in\mathcal G_W^+}\Xi_g$.
    An occurrence with no participants retains its aggregate operation
    and requests no monomial label. The bijection includes empty label
    sets. These are the choices used in
    \cite[proof of Theorem~5.2]{OpenAI2026PolynomialPEPS},
    \texttt{04-compression.tex}, lines 351--381 and 565--573.
\end{definition}

\begin{theorem}[Coefficients of chronological choices]
    \label{thm:peps_chronological_coefficients}
    \lean{TNLean.PEPS.PairEffect.SourceCircuit.coefficient_all_eq_prod_nonemptyGateLocations}
    \leanok
    \uses{def:peps_nonempty_gate_labels}
    For an all-party choice $\xi$, its coefficient is
    $c(\xi)=\prod_{g\in\mathcal G_W^+}c_g(\xi_g)$.
\end{theorem}

\begin{proof}\leanok
    \uses{def:peps_nonempty_gate_labels}
    Induction on the chronology gives the product of coefficients at
    exactly the selected occurrences. Unselected occurrences contribute
    one to this product. Restrict the product to $\mathcal G_W^+$.
\end{proof}

\begin{theorem}[Incident dependence of the actual residual]
    \label{thm:peps_incident_residual}
    \lean{TNLean.PEPS.PairEffect.SourceCircuit.restrict_partialResidual_eval_eq_of_projectAllChoices_eq}
    \leanok
    Hoist the sources in each chosen monomial into the fixed common
    source registers. Retain the auxiliary scalar owner together with
    the original party owners. For any selected set of owners, the
    restriction of the resulting actual residual depends only on the
    monomial labels at occurrences meeting that set.
    In particular, two all-party choices with the same incident labels
    give the same restricted linear map. No cancellation of source
    preparations is required.
\end{theorem}

\begin{proof}\leanok
    Induct on the chronology. At an incident gate the prescribed label
    identifies the chosen branch. A gate disjoint from the selected
    owners acts as the identity on their spectator registers, even
    before its sources are evaluated. An occurrence with no
    participants is independent of the supplied choice. Composition
    and the addition of spectator registers preserve the equality.
    Transport both maps to the fixed common source coordinates.
\end{proof}

\begin{theorem}[Tensor factorization after the original private preparation]
    \label{thm:peps_actual_private_factorization}
    \lean{TNLean.PEPS.PairEffect.Word.groupByPartyIso_eval_eq_tensorRestrictions,
      TNLean.PEPS.PairEffect.SourceCircuit.groupByPartyIso_partialSourceWord_eq_tensorPartialSourceLocalMaps}
    \leanok
    \uses{thm:peps_incident_residual}
    Let $\psi$ be a recorded product input and $\xi$ an all-party
    monomial choice. Prepare $\psi$ locally beside the fixed source
    registers, then apply the actual residual of $\xi$.
    Let $F_\xi$ be this source-to-output linear map, and let
    $U_{\rm in},U_{\rm out}$ collect all registers by owner in one
    exhaustive order without repetitions. Then
    \begin{align}
        U_{\rm out}F_\xi
        &=\left(\bigotimes_q F_{q,\xi|_q}\right)U_{\rm in}.
        \label{eq:peps_actual_local_factorization}
    \end{align}
    Each $F_{q,\xi|_q}$ is the actual singleton restriction after its
    local input preparation. The auxiliary scalar owner is retained
    as one factor. Empty original party sets and empty monomial label
    sets are allowed in this identity.
\end{theorem}

\begin{proof}\leanok
    \uses{thm:peps_incident_residual}
    A chronology without source preparations factors into its literal
    singleton restrictions. To prove this, separate the first owner
    from all remaining registers, factor the two restricted maps, and
    induct on the owner order. Apply this identity to the composition
    of the private preparation and the actual residual. Their
    singleton restrictions compose; incident dependence identifies
    the residual restriction with the constructed local map.
\end{proof}

% These entries establish exact chronological local factors. The sampled
% endpoint operators and their final physical-network contraction are separate.
